\begin{abstract}
\noindent
We develop the Semantic Mixture of Experts (SMoE), a white-box model of the conditional
distribution of daily returns in which linear-Gaussian experts with log-linear variances are weighted by
a multinomial-logit gate in observable state variables, and each expert is restricted to a
researcher-specified semantic group of covariates. We establish identifiability under the semantic
exclusions, geometric ergodicity of the canonical SMoE autoregression, a monotone ECM algorithm with
exact scores, and QMLE sandwich inference; every predictive functional has a closed form (mean, variance split into
within-regime and regime-disagreement risk, probability of a rise, VaR, expected shortfall and CRPS), and every forecast an exact additive attribution. We also prove that MSE-trained
price-level forecasters must lag: their maximal RMSE gain over the random walk is of order
$s^2/2$ for daily drift-to-noise ratio $s$, i.e.\ well below one percent. Projecting black-box density
forecasters onto the SMoE by KL minimisation yields a \emph{yardstick}: a fidelity ratio, with an
exact information-theoretic meaning, of how much of a black box's out-of-sample gain three readable
regimes reproduce. On nineteen equity, commodity, bond and crypto series (2005--2026) evaluated under a
leak-tested walk-forward protocol, the SMoE has the best average log-score rank among thirteen
forecasters, three of them zero-shot foundation models (two sizes of Chronos-Bolt and TimesFM~2.5),
and well-calibrated tails (1\% VaR passes the Kupiec test on 17 of 19 series), but no model forecasts
daily means or directions. A deep heterogeneous mixture-of-experts network is 92\% white by the
yardstick, and white-box projections of all three foundation models outperform the models themselves;
a release-date test finds no sign that the foundation models had memorised the evaluation period.
The white box shows that volatility, not momentum or drawdown, drives regime weights in the CSI~300
and that the short-horizon expert exhibits continuation in Chinese and reversal in developed markets.
Comparing fourteen specifications, three experts forecast best out of sample, although BIC prefers
two and a bootstrapped likelihood ratio four; a score test computed from the baseline fit shows that
the VIX belongs in the gate of all six international equity indices and most commodity futures but of
only one of the four A-share indices. With
5-minute data, a measurement-error result predicts the attenuation of the range-based daily variance
loading (predicted ratio $0.39$, estimated $0.40$); realised measures improve the SMoE, but a
single-regime Realized GARCH forecasts best. An adaptive conformal layer guarantees long-run VaR
coverage. Trading the CSI~300 ETF on the SMoE
predictive distribution cuts the maximum drawdown of buy-and-hold from $-46\%$ to $-26\%$ (index
futures: from $-45\%$ to $-38\%$), but no excess return survives deflation for the strategies tried.
In a lockbox period from July 2024 to October 2026, evaluated once after all choices were frozen, the
mixture models again lead on log score, with the SMoE second of thirteen behind LSTAR (average rank

$3.11$ against $3.05$); GARCH-$t$ matches its tail calibration (1\% VaR passes the Kupiec test on 18
of 19 series for both); and there is still no reliable trading alpha.

\medskip\noindent\textbf{Keywords:} mixture of experts; density forecasting; interpretability;
regime switching; conformal prediction; foundation models; CSI~300.

\noindent\textbf{JEL:} C22, C53, C58, G17.
\end{abstract}
